\chapter{ВВЕДЕНИЕ}
\label{chap:introduction}

Нелинейные социально-экономические процессы часто анализируются при помощи линейных регрессионных моделей, поскольку они дают интерпретируемые коэффициенты, позволяют сопоставлять факторы и обеспечивают стандартный набор диагностических процедур. Вместе с тем в задачах, где динамика режима меняется во времени, присутствуют циклические участки, локальные шоки и короткие окна наблюдений, сама по себе высокая точность подгонки не гарантирует корректность содержательного вывода. В этом смысле работа опирается одновременно на классические представления о нелинейной динамике, на стандартный регрессионный инструментарий и на литературу по эконометрике временных рядов \cite{strogatz,montgomery,hamilton,box-jenkins}.

Проблема исследования состоит в следующем: может ли линейный регрессионный анализ использоваться не как универсальная прогностическая модель, а как инструмент структурной диагностики нелинейного процесса. Для ответа на этот вопрос в работе сначала рассматриваются синтетические траектории с заранее заданным законом динамики, а затем построенный маршрут анализа переносится на реальные данные по выручке российских компаний. Такой переход позволяет разделить две задачи: в контролируемых условиях проверить границы интерпретации, а на эмпирике показать, как тот же подход работает там, где точное восстановление истинной модели невозможно \cite{strogatz,hamilton,wooldridge}.

Цель работы состоит в том, чтобы определить условия, при которых линейная регрессия сохраняет содержательную ценность для анализа нелинейных процессов, и показать, как эти условия переносятся на реальные социально-экономические данные.

Для достижения этой цели решаются следующие задачи:
\begin{enumerate}
\item построить единый вычислительный эксперимент для трёх классов нелинейных моделей;
\item сопоставить качество подгонки, структурную идентификацию и прогностическую устойчивость регрессии в разных динамических режимах;
\item исследовать границы интерпретируемости коротких окон и роль режимной диагностики;
\item перенести полученный маршрут анализа на реальные данные по выручке групп и отраслей;
\item оценить практическую значимость подхода как процедуры отбраковки неинтерпретируемых окон и выбора допустимого контура анализа.
\end{enumerate}

Объектом исследования являются нелинейные динамические процессы в синтетической и реальной социально-экономической среде. Предмет исследования --- условия применимости линейного регрессионного анализа для структурной диагностики таких процессов.

Гипотеза работы формулируется следующим образом: линейный регрессионный анализ может быть полезен для структурной идентификации нелинейных процессов, но его интерпретация и прогностическая применимость зависят от динамического режима, длины окна, локальной дисперсии и структуры лаговой зависимости; поэтому регрессия должна использоваться как часть режимно-зависимой процедуры анализа, а не как универсальный инструмент.

Научная новизна работы состоит в том, что регрессионный анализ рассматривается не как самодостаточная модель, а как элемент последовательной диагностической процедуры: сначала определяется режим процесса, затем проверяется интерпретируемость окна, после чего выбирается допустимый контур чтения коэффициентов. Практическая значимость работы связана с тем, что такая процедура снижает риск ложных выводов по коротким и неоднородным экономическим рядам \cite{montgomery,wooldridge}.

Структура работы соответствует поставленным задачам. Во второй главе приводится вычислительный эксперимент на синтетических траекториях. В третьей главе построенный подход переносится на реальные данные по выручке российских компаний. В заключении формулируются основные результаты и ограничения исследования.
